\documentclass[aps,prb,onecolumn,nofootinbib,superscriptaddress]{revtex4-2}
\usepackage{amsmath,amssymb,amsthm,mathtools,bm}
\usepackage{booktabs}
\usepackage[colorlinks=true,linkcolor=blue,citecolor=blue,urlcolor=blue]{hyperref}
\usepackage{microtype}
\newcommand{\Id}{\mathbf{1}}
\newcommand{\One}{\hat{\mathbf{1}}}
\newcommand{\Tr}{\operatorname{Tr}}
\newcommand{\tr}{\operatorname{tr}}
\newcommand{\dd}{\mathrm{d}}
\newcommand{\ii}{\mathrm{i}}
\newcommand{\e}{\mathrm{e}}
\newcommand{\Var}{\operatorname{Var}}
\newcommand{\ket}[1]{\lvert #1\rangle}
\newcommand{\bra}[1]{\langle #1\rvert}
\hypersetup{pdftitle={U(1) Kac-Moody Level and Renyi Entanglement in Monitored Gaussian States},
pdfauthor={Class-A Fermionic Adaptive-Circuit Project}}
\begin{document}
\title{\texorpdfstring{$U(1)$ Kac--Moody Level and R\'enyi Entanglement\\
in Monitored Gaussian States}{U(1) Kac-Moody Level and Renyi Entanglement in Monitored Gaussian States}}
\author{Class-A Fermionic Adaptive-Circuit Project}
\noaffiliation
\date{October 6, 2026}
\begin{abstract}
We derive the relation between interval entropies, charge fluctuations, and
the infrared data of a fermionic boundary theory. We first obtain exact
Gaussian-state formulas from the restricted correlation matrix, keeping
conditional observables distinct from those of a measurement-record mixture.
We specify chiral and nonchiral fermion Hamiltonians and actions, explain
their bosonization, and compare ordinary and chiral Luttinger liquids.
We then explain the replica construction of the interval entropy and derive
the cylinder charge variance from a normalized current algebra. We introduce
the full affine Kac--Moody algebra, its representations, and its relation
to the Virasoro algebra before specializing to physical charge. These
calculations identify different quantities: entropy measures the sum of
left- and right-moving central charges, while charge fluctuations measure
the current levels. Their relation for a charge-one complex fermion is
explained using the Sugawara construction and the common susceptibility
spectrum. An interface-defect example explains how logarithmic entanglement
can survive with a continuously varying coefficient and a gapped modular
spectrum, without changing the bulk central charge or current algebra.
Finally, we discuss one-wall versus two-wall normalization and
the interpretation of spatial entanglement contours. The Gaussian identities
are exact; their conformal interpretation requires additional assumptions
about the long-distance state.
\end{abstract}
\maketitle
\setcounter{tocdepth}{2}
\tableofcontents

\section{Gaussian reduced states and interval entropies}
\label{sec:gaussian}
The purpose of this note is to connect two descriptions of the same
subsystem. Microscopically, a number-conserving Gaussian state is specified
by its two-point function. At long distances, a critical interface may
instead be described by a conformal field theory (CFT). We will derive the
microscopic entropy and charge observables first, and only then introduce
the assumptions under which their logarithmic coefficients become CFT data.
We use the correlation-matrix convention of Bhuiyan, Pan, and Jian
(BPJ) \cite{BhuiyanPanJian2026}.

Three distinctions guide the argument. The eigenvalues of a reduced
density matrix are properties of a particular state and partition;
the central charge and current level are coefficients in local continuum
operator algebras; and a coefficient extracted from an interval entropy
or variance becomes one of those algebraic quantities only after the
geometry and infrared state have been identified. In particular,
``logarithmic'' does not by itself mean ``the homogeneous CFT vacuum.''
We keep the exact Gaussian calculation separate from that interpretation,
and later use a solvable defect to see why the separation matters.
All logarithms are natural. Ordinary single kets denote states throughout.

\subsection{Correlation convention and natural orbitals}
Let $\boldsymbol m$ denote a complete measurement record, and
$\hat\rho_{\boldsymbol m}$ its normalized conditioned Gaussian state. For
complex fermions obeying
$\{\hat\psi_i^\dagger,\hat\psi_j\}=\delta_{ij}$ and
$\{\hat\psi_i,\hat\psi_j\}=0$, define
\begin{equation}
 (G_{\boldsymbol m})_{ij}
 =\Tr(\hat\rho_{\boldsymbol m}\hat\psi_i^\dagger\hat\psi_j).
 \label{eq:G}
\end{equation}
Restricting $i,j$ to a region $A$ gives $G_{\boldsymbol m,A}$.
Here $A$ always denotes a subsystem, not an ancilla. Positivity and the
canonical anticommutation relations imply
$0\leq G_{\boldsymbol m,A}\leq\Id_A$. For example,
$v^\dagger Gv$ is the expectation of
$(\sum_i v_i\hat\psi_i)^\dagger(\sum_j v_j\hat\psi_j)$, and the
anticommutator gives the complementary upper bound.

Temporarily suppress the record label and diagonalize
\begin{equation}
 G_A=U\operatorname{diag}(\nu_a)U^\dagger,\qquad
 0\leq\nu_a\leq1,\qquad
 \hat f_a=\sum_i U_{ia}\hat\psi_i .
 \label{eq:natural-orbitals}
\end{equation}
Then
$\langle\hat f_a^\dagger\hat f_b\rangle
=(U^\dagger G_AU)_{ab}=\nu_a\delta_{ab}$.
The placement of conjugates in this formula is important. With the
definition in Eq.~\eqref{eq:G}, the occupation matrix used in the usual
one-body Hamiltonian ordering is $G_A^{\mathsf T}$.

Partial trace preserves Gaussianity. In the natural-orbital basis the
reduced state is a product of independent occupied/empty modes
\cite{Peschel2003}:
\begin{align}
 \hat\rho_A
 &=\bigotimes_a\left[
 (1-\nu_a)\ket{0_a}\bra{0_a}
 +\nu_a\ket{1_a}\bra{1_a}\right],
 \label{eq:rho-product}\\
 &=Z_A^{-1}\exp\!\left[-\sum_a\epsilon_a\hat f_a^\dagger\hat f_a\right],
 \qquad
 \epsilon_a=\log\frac{1-\nu_a}{\nu_a},\qquad
 Z_A=\prod_a(1+\e^{-\epsilon_a}).
 \label{eq:rho-modular}
\end{align}
One way to see the factorization is to use Wick's theorem: in a basis with
diagonal two-point functions, all occupation moments factor into products
of $\nu_a$. Gaussianity excludes independent higher connected correlations,
so these moments fix the full product density matrix. Diagonalizing a
two-point function alone would not justify this conclusion for a
non-Gaussian state.

In the original site basis write
$\hat\rho_A\propto\exp[-\sum_{ij}(h_A)_{ij}
\hat\psi_i^\dagger\hat\psi_j]$. Substitution of
Eq.~\eqref{eq:natural-orbitals} gives
\begin{equation}
 h_A=U^*\operatorname{diag}(\epsilon_a)U^{\mathsf T}
 =\left\{\log[(\Id_A-G_A)G_A^{-1}]\right\}^{\mathsf T},
 \qquad
 G_A=(\Id_A+\e^{h_A^{\mathsf T}})^{-1}.
 \label{eq:transpose}
\end{equation}
Thus the transpose matters for eigenvectors and real-space modular
evolution, although it does not change any spectral entropy below.
Eigenvalues $\nu_a=0,1$ are understood by continuity: the corresponding
mode is frozen and $\epsilon_a$ is infinite, not a finite matrix logarithm.

It is useful to distinguish the single-particle entanglement energies
$\epsilon_a$ from the many-body modular spectrum. Defining
$\hat{\mathcal H}_A=-\log\hat\rho_A$, its eigenvalues are
\begin{equation}
 \mathcal E_{\{n_a\}}=\log Z_A+\sum_a\epsilon_a n_a,
 \qquad n_a\in\{0,1\}.
 \label{eq:manybody-modular}
\end{equation}
A negative $\epsilon_a$ simply means that occupation is more likely than
vacancy; the probabilities $\exp(-\mathcal E_{\{n_a\}})$ remain between
zero and one. This modular operator is generally not the physical
Hamiltonian restricted to $A$, and its dimensionless energies do not
measure a physical excitation gap. The distinction is developed in
Ref.~\cite{PeschelEisler2009}, Secs.~3--5, and will be important for the
defect example in Sec.~\ref{sec:defect}.

\subsection{Taking the trace of a power}
The many-body eigenvalues of Eq.~\eqref{eq:rho-product} are
$\prod_a\nu_a^{n_a}(1-\nu_a)^{1-n_a}$ with $n_a=0,1$. Consequently,
\begin{align}
 \Tr\hat\rho_A^q
 &=\prod_a[\nu_a^q+(1-\nu_a)^q],\\
 S_q(A)&=\frac{1}{1-q}\sum_a
       \log[\nu_a^q+(1-\nu_a)^q],\qquad q>0,\quad q\ne1 .
 \label{eq:renyi}
\end{align}
Here $q$ is always the R\'enyi index, not electric charge or a band label.
Differentiating the logarithm at $q=1$ gives
\begin{equation}
 S_1(A)=\sum_a h_1(\nu_a),\qquad
 h_1(\nu)=-\nu\log\nu-(1-\nu)\log(1-\nu),
 \label{eq:von-neumann}
\end{equation}
with $0\log0=0$. For a single mode at $\nu=1/2$, all R\'enyi entropies
equal $\log2$; a frozen mode contributes zero. These are exact finite-system
identities, independent of criticality.

For a pure global state, the reduced-state eigenvalues are the squared
Schmidt coefficients. The simplest example is the one-particle state
\begin{equation}
 \ket{\Psi_\nu}=\sqrt\nu\ket{1_A0_{\bar A}}
             +\sqrt{1-\nu}\ket{0_A1_{\bar A}}.
 \label{eq:schmidt-pair}
\end{equation}
Tracing out $\bar A$ gives a single factor of
Eq.~\eqref{eq:rho-product}. The same uncertain location of the particle
produces both $h_1(\nu)$ entanglement and $\nu(1-\nu)$ charge variance.
Many Gaussian modes repeat this construction in suitable bases, but
their entropy and variance add with different weights.

Increasing $q$ emphasizes the largest many-body eigenvalues $p_\alpha$.
In particular,
\begin{equation}
 S_\infty=-\log p_{\max}
 =-\sum_a\log\max(\nu_a,1-\nu_a)
 =\sum_a\log(1+\e^{-|\epsilon_a|}).
 \label{eq:single-copy}
\end{equation}
By contrast $S_0=\log\operatorname{rank}\hat\rho_A$ counts every
nonzero eigenvalue equally. In a finite Gaussian system it is $\log2$
times the number of modes with $0<\nu_a<1$. Its limit need not commute
with a continuum or large-size limit: an arbitrarily small but nonzero
Schmidt weight still contributes to $S_0$.

If the global conditioned state is pure, these reduced-state entropies are
bipartite entanglement entropies. If it is mixed, the same spectral formulas
are correct, but include global mixedness. For example,
$G_A=\Id_A/2$ gives $S_q=|A|\log2$ even for an entirely uncorrelated
maximally mixed state. A logarithmic entanglement interpretation therefore
needs a separate physical argument about the state, not just a Gaussian
formula.

\section{Charge statistics and measurement-record averages}
\label{sec:fcs}
\subsection{Generating function and cumulants}
In units where a physical fermion has charge one, let
$\hat Q_A=\sum_{i\in A}\hat\psi_i^\dagger\hat\psi_i
=\sum_a\hat f_a^\dagger\hat f_a$. Each natural orbital contributes a
Bernoulli occupation. Its characteristic function is therefore
\begin{align}
 \mathcal Z_A(\lambda)
 &=\Tr(\hat\rho_A\e^{\ii\lambda\hat Q_A})
 =\prod_a(1-\nu_a+\e^{\ii\lambda}\nu_a)\\
 &=\det[\Id_A-G_A+\e^{\ii\lambda}G_A].
 \label{eq:fcs}
\end{align}
This is the Gaussian full-counting-statistics determinant; Klich's trace
identity supplies a general operator derivation \cite{Klich2003}.
Our product-state derivation also shows why a single-particle determinant
can encode a many-body charge distribution.

Define the cumulants by
$\kappa_j=(-\ii\partial_\lambda)^j\log\mathcal Z_A|_{\lambda=0}$.
For one orbital, put $t=\ii\lambda$ and
$g(t)=\log(1-\nu+\nu\e^t)$. If
$p(t)=\nu\e^t/(1-\nu+\nu\e^t)$, then
$g'=p$ and $p'=p(1-p)$. Repeated differentiation at $t=0$ gives
\begin{align}
 \kappa_1&=\tr G_A,&
 F_A\equiv\kappa_2&=\tr[G_A(\Id_A-G_A)],\nonumber\\
 \kappa_3&=\tr[G_A(\Id_A-G_A)(\Id_A-2G_A)],&
 \kappa_4&=\tr[G_A(\Id_A-G_A)(\Id_A-6G_A+6G_A^2)].
 \label{eq:cumulants}
\end{align}
In particular $F_A=\sum_a\nu_a(1-\nu_a)$ is the quantum variance of
charge in this conditioned state. Relations between free-fermion
entanglement and charge statistics are developed in
Refs.~\cite{KlichLevitov2009,SongFlindtRachel2011,SongRachelFlindtKlichLaflorencieLeHur2012}.

\subsection{Which average is being taken?}
Restore the record label and let an overline denote the Born average:
$\overline{X}=\sum_{\boldsymbol m}p_{\boldsymbol m}X_{\boldsymbol m}$.
The mean conditional entropy, annealed replica moment, and entropy of
the Born mixture are three different quantities:
\begin{align}
 \overline{S_q(A)}
 &=\frac{1}{1-q}\,
   \overline{\log\Tr\hat\rho_{\boldsymbol m,A}^{\,q}},
 \label{eq:conditional-average}\\
 S_q^{\mathrm{ann}}(A)
 &=\frac{1}{1-q}\log
   \overline{\Tr\hat\rho_{\boldsymbol m,A}^{\,q}},\\
 S_q^{\mathrm{mix}}(A)
 &=\frac{1}{1-q}\log\Tr
   \left(\overline{\hat\rho_{\boldsymbol m,A}}\right)^q.
 \label{eq:mixture}
\end{align}
For $q>1$, Jensen's inequality gives
$\overline{S_q}\geq S_q^{\mathrm{ann}}$; the mixture introduces cross terms
between different states and is a distinct construction. For example, an
equal ensemble of empty and occupied one-mode pure states has zero
conditional entropy but mixture entropy $\log2$.

A mixture of Gaussian states is generally non-Gaussian. Thus inserting
$\overline G$ into Eq.~\eqref{eq:renyi} computes the entropy of the
Gaussian state with that covariance, not generally that of the actual
mixture. An equal mixture of the two-mode states $\ket{00}$ and
$\ket{11}$ has $\overline G=\Id/2$: the actual entropy is $\log2$,
whereas the Gaussian formula gives $2\log2$.

Charge fluctuations separate just as explicitly. Write
$\mu_{\boldsymbol m}=\tr G_{\boldsymbol m,A}$ and
$F_{\boldsymbol m,A}=\langle\hat Q_A^2\rangle_{\boldsymbol m}
-\mu_{\boldsymbol m}^2$. The law of total variance gives
\begin{equation}
 \Var_{\mathrm{mix}}(Q_A)
 =\overline{F_A}
   +\underbrace{\overline{\mu_{\boldsymbol m}^2}
                    -\overline{\mu_{\boldsymbol m}}^{\,2}}
               _{\text{record-to-record charge wandering}}.
 \label{eq:total-variance}
\end{equation}
The current-algebra comparison below uses the first term.
Averaging translations of a strip within one trajectory is another operation;
when needed we denote it explicitly by
$N_y^{-1}\sum_{y_0}$, reserving the overline for records.

\section{Continuum fermions, bosonization, and Luttinger liquids}
\label{sec:continuum-models}
Before extracting continuum data from an entropy, it is useful to specify
the theories that could carry those data. We use $y$ for the coordinate
along an interface and set $\hbar=1$. Continuum fields and their mode
operators are written without hats in the standard field-theory notation;
$\hat H$ still denotes a many-body Hamiltonian. Path-integral fermion
fields are Grassmann variables. None of the Hamiltonians in this section
is being identified with the modular Hamiltonian
$\hat{\mathcal H}_A$ or with a generator of the adaptive circuit.

\subsection{A Dirac fermion and a single chiral fermion}
Linearizing a one-dimensional spinless band about its two Fermi points
gives $c(y)\simeq\e^{\ii k_Fy}\psi_+(y)+\e^{-\ii k_Fy}\psi_-(y)$.
The slowly varying fields obey
$\{\psi_\sigma(y),\psi_{\sigma'}^\dagger(y')\}
=\delta_{\sigma\sigma'}\delta(y-y')$. We use the Weyl basis with
$\psi_+=\psi_R$,
$\psi_-=\psi_L$, and $\Psi=(\psi_+,\psi_-)^{\mathsf T}$.
For equal positive velocities $v_F$, the normal-ordered Hamiltonian is
\begin{align}
 \hat H_{\mathrm D}
 &=-\ii v_F\int\dd y\,
       \left(:\!\psi_+^\dagger\partial_y\psi_+\!:
             -:\!\psi_-^\dagger\partial_y\psi_-\!:\right)
 =\int\dd y:\!\Psi^\dagger(-\ii v_F\sigma^3\partial_y)\Psi\!:,
 \label{eq:dirac-hamiltonian}\\
 \mathcal S_{\mathrm D}
 &=\int\dd t\,\dd y\,
    \left[\ii\psi_+^\dagger(\partial_t+v_F\partial_y)\psi_+
         +\ii\psi_-^\dagger(\partial_t-v_F\partial_y)\psi_-\right],
 \label{eq:dirac-action}\\
 \mathcal S_{\mathrm D,E}
 &=\int\dd\tau\,\dd y\,
    \left[\psi_+^\dagger(\partial_\tau-\ii v_F\partial_y)\psi_+
         +\psi_-^\dagger(\partial_\tau+\ii v_F\partial_y)\psi_-\right].
 \label{eq:dirac-euclidean}
\end{align}
Normal ordering subtracts the filled reference sea; it is essential for
the continuum density fluctuations. The equations of motion give
$\psi_+(y,t)=\psi_+(y-v_Ft,0)$ and
$\psi_-(y,t)=\psi_-(y+v_Ft,0)$. Thus the signs in the Hamiltonian
specify actual propagation, not merely names for two fields
\cite{Senechal1999}, Sec.~3.1, Eq.~(19).

A purely right-moving complex fermion retains only the $+$ term:
\begin{equation}
 \hat H_R=-\ii v_R\int\dd y:\!\psi_R^\dagger\partial_y\psi_R\!:,
 \qquad
 \mathcal S_R=\int\dd t\,\dd y\,
               \ii\psi_R^\dagger(\partial_t+v_R\partial_y)\psi_R.
 \label{eq:chiral-fermion}
\end{equation}
Its Euclidean action is the corresponding single term in
Eq.~\eqref{eq:dirac-euclidean}. A single complex chiral fermion has
$(c_R,c_L)=(1,0)$; the Dirac theory has $(1,1)$.
The customary statement ``a Dirac fermion has $c=1$'' uses the
nonchiral convention $c_R=c_L=c$, not $c_R+c_L=1$.
Likewise the physical charge-one current has level one in each occupied
chiral sector. The algebraic meanings of $c$ and $k$ are derived below.

An isolated chiral theory is an effective edge theory with a microscopic
regulator supplied, for example, by a two-dimensional bulk. It is not a
standalone one-dimensional lattice band obtained by forgetting its
opposite Fermi point. On a strip the opposite edge can supply the
counterpropagating branch, without putting both branches at the same
spatial interface. This distinction is precisely why local wall
chirality can coexist with zero net chirality for the two-wall system.

\subsection{Bosonization conventions and the fermion dictionary}
We fix the bosonization normalization explicitly before using it.
For a real scalar $\Phi$ with canonical momentum $\Pi$, introduce its
dual $\varphi$ and chiral components by
\begin{equation}
 [\Phi(y),\Pi(y')]=\ii\delta(y-y'),\qquad
 \partial_y\varphi=\Pi,\qquad
 \Phi_\pm=\frac{\Phi\mp\varphi}{2},\qquad
 \Phi=\Phi_++\Phi_-.
 \label{eq:user-boson-fields}
\end{equation}
On a circle the dual includes winding and a zero-mode convention; the
indefinite integral of $\Pi$ specifies its local derivative, not all
global sectors. On the line our chiral commutators are
\begin{equation}
 [\Phi_\pm(y),\Phi_\pm(y')]
      =\pm\frac{\ii}{4}\operatorname{sgn}(y-y'),\qquad
 [\Phi_+(y),\Phi_-(y')]=\frac{\ii}{4}.
 \label{eq:user-boson-commutators}
\end{equation}
The constant cross-commutator fixes the relative fermionic sign. With
short-distance cutoff $a$, the dictionary is
\begin{equation}
 \psi_\pm(y)=\frac{1}{\sqrt{2\pi a}}
       :\!\e^{\pm\ii\sqrt{4\pi}\Phi_\pm(y)}\!:,
 \qquad
 \rho_\pm(y)=:\!\psi_\pm^\dagger\psi_\pm\!:
       =\frac{1}{\sqrt\pi}\partial_y\Phi_\pm(y).
\label{eq:user-bosonization}
\end{equation}
Here $a$ denotes the ultraviolet length cutoff, reserving $\epsilon_a$
for modular energies; it has the role of the cutoff $\epsilon$ in the
bosonization cheat sheet.
For example, the commutator of the two vertex exponents is $\ii\pi$,
so exchanging them supplies a minus sign. An alternative convention
takes commuting cross-chiral bosons and introduces explicit Klein
factors to supply this sign. We do not add those factors on top of
Eq.~\eqref{eq:user-boson-commutators}; that would count the same sign
twice. Point splitting and normal ordering give the density identity in
Eq.~\eqref{eq:user-bosonization}; see
Ref.~\cite{Senechal1999}, Secs.~4--5, for the operator construction.

At the free point these relations give
\begin{align}
 \hat H_{\mathrm B}
 &=\frac{v_F}{2}\int\dd y\,[\Pi^2+(\partial_y\Phi)^2]
 =v_F\int\dd y\,[(\partial_y\Phi_+)^2+(\partial_y\Phi_-)^2],
 \label{eq:free-boson-hamiltonian}\\
 \mathcal S_{\mathrm B}
 &=\frac12\int\dd t\,\dd y\,
       [v_F^{-1}(\partial_t\Phi)^2-v_F(\partial_y\Phi)^2],
 \label{eq:free-boson-action}\\
 \mathcal S_{\mathrm B,E}
 &=\frac12\int\dd\tau\,\dd y\,
       [v_F^{-1}(\partial_\tau\Phi)^2+v_F(\partial_y\Phi)^2].
 \label{eq:free-boson-euclidean}
\end{align}
One obtains the action by integrating $\Pi$ out of
$\int\dd t\,\dd y\,[\Pi\partial_t\Phi-\mathcal H]$.
The conserved physical density and current are
\begin{equation}
 \rho=\rho_++\rho_-=\frac{\partial_y\Phi}{\sqrt\pi},\qquad
 j=v_F(\rho_+-\rho_-)=-\frac{v_F\Pi}{\sqrt\pi},\qquad
 \partial_t\rho+\partial_yj=0.
 \label{eq:user-charge-current}
\end{equation}
The plus sign in the density and the factors of $\sqrt\pi$ fix our
convention throughout; changing either would change the charge dictionary.

The regulated equal-time contractions are
\begin{equation}
 G_\pm(y)=\frac{1}{4\pi}\log\frac{a}{a\mp\ii y},\qquad
 G(y)=\frac{1}{4\pi}\log\frac{a^2}{a^2+y^2}.
 \label{eq:user-boson-propagators}
\end{equation}
These expressions fix an additive constant, equivalently a normal-ordering
scale; physical neutral vertex correlators do not depend on the otherwise
infrared-divergent scalar zero mode. Gaussian contraction gives
$:\!\e^A\!:\!\e^B\!:=:\!\e^{A+B}\!:\e^{\langle AB\rangle}$
for linear bosonic exponents in a common radial ordering. Hence a chiral
vertex $:\!\e^{\ii\sqrt{4\pi}\alpha\Phi_+}\!:$ has weight
$h=\alpha^2/2$, while the full vertex
$:\!\e^{\ii\sqrt{4\pi}\alpha\Phi}\!:$ has
$h=\bar h=\alpha^2/2$ and scaling dimension $\Delta=\alpha^2$.
The fermion vertex at $\alpha=1$ has chiral weight $1/2$, as required.
We define these conformal terms more formally in Sec.~\ref{sec:replica}.

\subsection{The ordinary Luttinger liquid}
Bosonization becomes particularly useful because density interactions,
quartic in fermions, are quadratic in the boson. Consider the gapless
density terms, with interbranch coupling $g_B$ and intrabranch coupling $g_F$:
\begin{equation}
 \hat H_{\mathrm{int}}=\int\dd y\,
       [2g_B:\!\rho_+\rho_-\!:
                +g_F:\!(\rho_+^2+\rho_-^2)\!:].
 \label{eq:density-interactions}
\end{equation}
The density dictionary yields
$2\rho_+\rho_-=[(\partial_y\Phi)^2-\Pi^2]/(2\pi)$ and
$\rho_+^2+\rho_-^2=[(\partial_y\Phi)^2+\Pi^2]/(2\pi)$.
Adding Eq.~\eqref{eq:free-boson-hamiltonian} therefore gives
\begin{align}
 \hat H_{\mathrm{LL}}
 &=\frac{u}{2}\int\dd y\,
       [K_{\mathrm{LL}}\Pi^2+K_{\mathrm{LL}}^{-1}(\partial_y\Phi)^2],
 \label{eq:ll-hamiltonian}\\
 uK_{\mathrm{LL}}&=v_F+\frac{g_F-g_B}{\pi},\qquad
 \frac{u}{K_{\mathrm{LL}}}=v_F+\frac{g_F+g_B}{\pi},
 \label{eq:ll-parameters}\\
 \mathcal S_{\mathrm{LL}}
 &=\frac{1}{2K_{\mathrm{LL}}}\int\dd t\,\dd y\,
       [u^{-1}(\partial_t\Phi)^2-u(\partial_y\Phi)^2],\nonumber\\
 \mathcal S_{\mathrm{LL},E}
 &=\frac{1}{2K_{\mathrm{LL}}}\int\dd\tau\,\dd y\,
       [u^{-1}(\partial_\tau\Phi)^2+u(\partial_y\Phi)^2].
 \label{eq:ll-actions}
\end{align}
Both stiffnesses must be positive. $K_{\mathrm{LL}}=1$ at the
noninteracting point; the density remains $\partial_y\Phi/\sqrt\pi$,
while continuity now fixes $j=-uK_{\mathrm{LL}}\Pi/\sqrt\pi$.
The field rescaling $\widetilde\Phi=\Phi/\sqrt{K_{\mathrm{LL}}}$,
$\widetilde\varphi=\sqrt{K_{\mathrm{LL}}}\varphi$ restores the free
scalar action but cannot remove $K_{\mathrm{LL}}$ from the physical
charge. This is the standard spinless Luttinger liquid
\cite{Senechal1999}, Sec.~6.1, Eqs.~(115)--(119).

The propagating eigenfields are
$\widetilde\Phi_\pm=(\widetilde\Phi\mp\widetilde\varphi)/2$.
Their normalized oscillator algebras have the free form, but the physical
density is $\sqrt{K_{\mathrm{LL}}/\pi}\,
\partial_y(\widetilde\Phi_++\widetilde\Phi_-)$.
Consequently the physical current levels are
$k_R=k_L=K_{\mathrm{LL}}$, whereas $c_R=c_L=1$.
There is also a direct derivation of the variance. For an interval,
\begin{equation}
 Q_A-\langle Q_A\rangle
   =\frac{\delta\Phi(\ell)-\delta\Phi(0)}{\sqrt\pi},\qquad
 F_A=\frac{K_{\mathrm{LL}}}{\pi^2}
          \log\frac{D_L(\ell)}{a}+\cdots,
 \quad D_L(\ell)=\frac{L}{\pi}\sin\frac{\pi\ell}{L}.
 \label{eq:ll-variance}
\end{equation}
Here $\delta\Phi=\Phi-\langle\Phi\rangle$. The action multiplies the scalar covariance in
Eq.~\eqref{eq:user-boson-propagators} by $K_{\mathrm{LL}}$; taking
the variance of the field difference and mapping to a circle gives the
last expression. The entropy retains the nonchiral $c=1$ coefficient.
Thus an ordinary Luttinger liquid supplies a concrete, nonchiral theory
whose charge-fluctuation coefficient varies continuously while its
central charge does not \cite{SongRachelFlindtKlichLaflorencieLeHur2012}.

Not every fermion interaction is included in this gapless quadratic
action. A right--left mixing bilinear becomes a vertex with the form
$\cos(\sqrt{4\pi}\Phi+\delta)$, and a commensurate spinless umklapp
term has the form $\cos(\sqrt{16\pi}\Phi+\delta')$, with scaling
dimensions $K_{\mathrm{LL}}$ and $4K_{\mathrm{LL}}$, respectively.
The phases depend on the relative fermion convention; they do not change
the dimensions. The spinless four-fermion term requires point splitting,
not a product of identical fields at exactly the same point. Such
perturbations may gap the liquid when allowed and relevant. A single
isolated chiral branch has no opposite mover with which to form this
ordinary Dirac mass. For distant opposite walls, a mixing term instead
requires interwall coupling.

\subsection{The chiral Luttinger liquid and its charge normalization}
The chiral theory has only one propagating density field. In the same
normalization as $\Phi_+$, a right-moving compact boson can be written
\begin{align}
 \mathcal S_{R,\mathsf M}
 &=-\mathsf M\int\dd t\,\dd y\,
       [\partial_t\Phi_R\partial_y\Phi_R+v_R(\partial_y\Phi_R)^2],
 \label{eq:chiral-boson-action}\\
 \hat H_{R,\mathsf M}&=v_R\mathsf M\int\dd y\,(\partial_y\Phi_R)^2,
 \qquad
 \mathcal S_{R,\mathsf M,E}=\mathsf M\int\dd\tau\,\dd y\,
       [\ii\partial_\tau\Phi_R\partial_y\Phi_R
                              +v_R(\partial_y\Phi_R)^2].
 \label{eq:chiral-boson-euclidean}
\end{align}
Here $\mathsf M>0$ is a stiffness together with a specified compact charge
lattice, not the Luttinger parameter $K_{\mathrm{LL}}$.
The negative sign of the real-time kinetic term is fixed by our choice
that a right mover satisfies $\partial_t\Phi_R=-v_R\partial_y\Phi_R$,
up to a removable spatially uniform mode. Reversing spatial orientation
reverses the kinetic sign, not the positive Hamiltonian. Quantizing the
first-order action gives
\begin{equation}
 [\Phi_R(y),\Phi_R(y')]=\frac{\ii}{4\mathsf M}
                         \operatorname{sgn}(y-y'),\qquad
 \rho_R=\frac{\partial_y\Phi_R}{\sqrt\pi},\qquad
 [\rho_R(y),\rho_R(y')]
       =-\frac{\ii}{2\pi\mathsf M}\partial_y\delta(y-y').
 \label{eq:chiral-boson-bracket}
\end{equation}
Fourier modes of this density obey the current algebra with
$k_R=1/\mathsf M$, while the oscillator sector has $c_R=1$.
At $\mathsf M=1$ this is precisely the bosonized charge-one chiral
complex fermion, not an additional interacting sector.

For comparison, a Laughlin edge at filling $1/\mathsf M$ with odd
integer $\mathsf M$ has an electron annihilator
$:\!\e^{\ii\mathsf M\sqrt{4\pi}\Phi_R}\!:$ of weight
$h=\mathsf M/2$ and charge $-1$. Its fundamental smaller vertex has
fractional charge. This illustrates why the scalar normalization,
compactification, and physical electron operator must be specified
together. It is an anomalous edge supported by a topological bulk,
not a standalone nonchiral compact-boson theory
\cite{Wen1995}, Sec.~3.1, and \cite{TongQHE}, Secs.~6.1.3--6.1.4,
with a sign reversal relative to conventions that orient the edge
oppositely. In the one-component topological $K$-matrix notation this
edge has $K_{\mathrm{edge}}=(\mathsf M)$: its quantized compactification
data are not the tunable $K_{\mathrm{LL}}$. For a fixed single Laughlin
branch, density interactions can renormalize $v_R$ without changing
$\mathsf M$. We are not assuming such fractional topological order for
the adaptive free-fermion protocol.

Both Luttinger descriptions are relevant to interpreting the observables,
but in different ways. The ordinary liquid is the natural counterexample
to equating a variance coefficient with chirality or with central charge.
The $\mathsf M=1$ chiral liquid is a useful candidate low-energy
description of an individual charge-one free-fermion wall. Neither
bosonization nor Gaussianity proves that a conditioned circuit state is
the vacuum of that local continuum theory. An interacting Luttinger
liquid is Gaussian in bosons, but generally not Gaussian in the original
fermions; the exact correlation-matrix entropy formulas of
Sec.~\ref{sec:gaussian} cannot be extended to it merely because its bosonic
action is quadratic.

\section{Replica derivation of the interval entropy}
\label{sec:replica}
The Gaussian identities do not require a continuum theory. We now assume
that the long-distance interval state is described by the vacuum of a
$1+1d$ CFT on a circle, or has the same leading interval scaling.
A generic monitored Gaussian state need not satisfy this assumption.
The replica method and its conformal implementation are reviewed in
Refs.~\cite{CalabreseCardy2004,CalabreseCardy2009}. To retain chirality
explicitly we follow the left/right separation discussed in
Ref.~\cite{CastroDetournayIqbalPerlmutter2014}, Sec.~II.

\subsection{The conformal ingredients}
We first describe the language needed for the replica calculation.
In two Euclidean dimensions a local conformal change of coordinate is
an analytic map $z\mapsto f(z)$ together with its antiholomorphic
counterpart. A primary field of weights $(h,\bar h)$ obeys
\begin{equation}
 \mathcal O(z,\bar z)=f'(z)^h\bar f'(\bar z)^{\bar h}
                \mathcal O(f(z),\bar f(\bar z)).
 \label{eq:primary-transformation}
\end{equation}
Its scaling dimension is $\Delta=h+\bar h$, while $h-\bar h$ is its
conformal spin. Translation, scaling, and special conformal invariance
then fix a conjugate pair's plane two-point function to
$C_{\mathcal O}(z_1-z_2)^{-2h}(\bar z_1-\bar z_2)^{-2\bar h}$.
The holomorphic and antiholomorphic dependencies continue to opposite
propagation directions in real time. Which is called right-moving is a
coordinate convention; neither label refers to the spatial position of a
domain wall.

The operator product expansion organizes what happens when two
insertions approach each other. The stress tensor generates local
coordinate changes, and its singular product with a primary is
\begin{equation}
 T(z)\mathcal O(u,\bar u)\sim
 \frac{h\mathcal O(u,\bar u)}{(z-u)^2}
 +\frac{\partial_u\mathcal O(u,\bar u)}{z-u}.
 \label{eq:T-primary}
\end{equation}
Integrating this expression around $u$ produces the infinitesimal
transformation of $\mathcal O$; applying the same contour to many insertions
gives the Ward identity used below. This is why a singular local product
can determine a global two-point scaling law. Radial quantization treats
concentric circles as equal-time slices and expands
\begin{equation}
 T(z)=\sum_{n\in\mathbb Z}L_nz^{-n-2},\qquad
 [L_m,L_n]=(m-n)L_{m+n}
       +\frac{c}{12}m(m^2-1)\delta_{m+n,0}\One.
 \label{eq:virasoro}
\end{equation}
The $L_n$ generate conformal transformations; the central term commutes
with all of them. There is a second commuting algebra for $\bar T$.
The central charge is consequently not a charge eigenvalue of the state.
It fixes the stress-tensor fluctuation strength and the anomalous term in
its transformation law. Standard derivations of these ingredients are
given in Ref.~\cite{Ginsparg1989}, Secs.~2--3, and
Ref.~\cite{DiFrancescoMathieuSenechal1997}.

\subsection{From gluing replicas to twist fields}
For an integer $q>1$, the Euclidean path integral for $\hat\rho_A$ is cut
along an interval $A=[u,v]$ at a fixed time. Multiplication of density
matrices glues the upper bank of one cut to the lower bank of the next;
the trace closes the sequence after $q$ sheets. Thus
\begin{equation}
 \Tr\hat\rho_A^q=\frac{Z[\mathcal R_q]}{Z[\mathcal R_1]^q},
 \label{eq:replica-gluing}
\end{equation}
where $\mathcal R_q$ is the branched surface. Equivalently, work on one
plane with $q$ copies of the CFT and insert a twist field $\mathcal T_q(u)$
and an inverse twist $\widetilde{\mathcal T}_q(v)$. Encircling one insertion
cycles the copy index. The ratio in Eq.~\eqref{eq:replica-gluing} is their
two-point function, with cutoff-dependent normalization.

To find the twist dimension for one holomorphic branch of central charge
$c$, uniformize the branched surface with
\begin{equation}
 f(z)=\left(\frac{z-u}{z-v}\right)^{1/q}.
 \label{eq:uniformization}
\end{equation}
A loop around $u$ changes sheets, while $f$ resolves the branch into a
single-valued local coordinate. The holomorphic stress tensor transforms as
\begin{equation}
 T(z)=f'(z)^2T(f(z))+\frac{c}{12}\{f,z\},\qquad
 \{f,z\}=\frac{f'''}{f'}-\frac32\left(\frac{f''}{f'}\right)^2.
 \label{eq:schwarzian}
\end{equation}
The plane vacuum has $\langle T(f)\rangle=0$. Since a fractional linear
map has zero Schwarzian and
$\{g^\alpha,g\}=(1-\alpha^2)/(2g^2)$, Eq.~\eqref{eq:uniformization} gives
\begin{equation}
 \{f,z\}=\frac{1-q^{-2}}{2}
 \frac{(u-v)^2}{(z-u)^2(z-v)^2}.
 \label{eq:schwarzian-result}
\end{equation}
The replica stress tensor sums over all $q$ sheets. Therefore
\begin{equation}
 \langle T_{\mathrm{tot}}(z)\rangle_{\mathcal R_q}
 =\frac{c}{24}\left(q-\frac1q\right)
   \frac{(u-v)^2}{(z-u)^2(z-v)^2}.
 \label{eq:replica-stress}
\end{equation}

For two primary insertions of weight $h$, the conformal Ward identity
expresses a stress-tensor insertion in terms of the two-point function:
\begin{align}
 \frac{\langle T_{\mathrm{tot}}(z)\mathcal T(u)
                      \widetilde{\mathcal T}(v)\rangle}
      {\langle\mathcal T(u)\widetilde{\mathcal T}(v)\rangle}
 &=\sum_{\zeta=u,v}\left[
    \frac{h}{(z-\zeta)^2}
    +\frac{1}{z-\zeta}\partial_\zeta
       \log (u-v)^{-2h}\right]\nonumber\\
 &=h\,\frac{(u-v)^2}{(z-u)^2(z-v)^2}.
 \label{eq:ward}
\end{align}
The derivatives contribute $-2h/(u-v)$ and $2h/(u-v)$, respectively;
combining the fractions gives the last line. Matching
Eq.~\eqref{eq:replica-stress} then determines
\begin{equation}
 h_q=\frac{c}{24}\left(q-\frac1q\right).
 \label{eq:twist-weight}
\end{equation}
The factor of $q$ from the replica stress tensor is essential. Missing it
would give the wrong R\'enyi coefficient. If both chiralities are present,
their weights are obtained separately with $c_R$ and $c_L$; their sum
sets the equal-time decay exponent.

\subsection{Plane-to-cylinder map and the R\'enyi coefficient}
On the plane, a single chiral twist pair has correlator proportional to
$(z_1-z_2)^{-2h_q}$. Let the circle circumference be $L=N_y$ in lattice
units, take $w=\tau+\ii y$ with Euclidean time rescaled by the branch
velocity, and map
\begin{equation}
 z=\exp(2\pi w/L),\qquad
 \mathcal T_q^{\mathrm{cyl}}(w)
 =\left(\frac{\dd z}{\dd w}\right)^{h_q}\mathcal T_q(z).
 \label{eq:cylinder-map}
\end{equation}
For equal-time endpoints at $y=0,\ell$,
$|z_1-z_2|=2\sin(\pi\ell/L)$ and
$|z'_1z'_2|^{h_q}=(2\pi/L)^{2h_q}$. Their ratio gives
\begin{equation}
 \Tr\hat\rho_{[0,\ell]}^q
 =C_q\left[\frac{L}{\pi a}\sin\left(\frac{\pi\ell}{L}\right)\right]^{-2h_q},
 \qquad
 X_L(\ell)=\log\left[\frac{L}{\pi a}
                          \sin\left(\frac{\pi\ell}{L}\right)\right].
 \label{eq:replica-cylinder}
\end{equation}
The microscopic cutoff $a$ and $C_q$ encode endpoint physics; the trace
uses a consistent regulator supplied by the microscopic system.
Taking $\log/(1-q)$ and using
$(q-1/q)/(q-1)=1+1/q$ gives, for a single branch,
\begin{equation}
 S_q^{(w)}(\ell)=\frac{c_w}{12}\left(1+\frac1q\right)X_L(\ell)
 +b_{w,q}+\cdots .
 \label{eq:one-branch-renyi}
\end{equation}
Analytic continuation from integer replicas to $q>0$ is the usual replica
assumption. In the limit $q\to1$ the coefficient is $c_w/6$.

The geometry is part of this result. There are two entangling endpoints
in an interval on a circle, even for one chiral branch. A nonchiral system
on a half-line, with an interval ending on a physical boundary, instead
uses a twist one-point function in the half-plane. Its logarithmic
coefficient is half that of a bulk interval, with an additional boundary
constant \cite{CalabreseCardy2009}, Sec.~3.2. A defect placed at an
entangling endpoint is a further change: the homogeneous two-point
function above no longer suffices. The half-chain defect calculation
below must therefore be compared with its own homogeneous half-chain,
not with the full-strip coefficient by a bare numerical identification.

For independent branches intersected by the region, leading contributions
add. Defining $c_\Sigma$ as the sum over both orientations and all walls,
\begin{equation}
 S_q(\ell)=a_{S_q}X_L(\ell)+b_q+\cdots,\qquad
 a_{S_q}=\frac{c_\Sigma}{12}\left(1+\frac1q\right).
 \label{eq:renyi-slope}
\end{equation}
Independence is an infrared assumption: interwall coupling or extra
gapless sectors changes the content being counted. Finite temperature,
an excited state, or a finite correlation length can also replace the
vacuum chord law. Gaussianity alone fixes none of these features.

In a monitored ensemble, the same coefficient passes to
$\overline{S_q}$ if the contributing records share it and their averaged
corrections remain subleading. A logarithmic ensemble mean alone does not
prove the same conformal description for every record. This use of
Eq.~\eqref{eq:renyi-slope} always concerns the conditional average of
Eq.~\eqref{eq:conditional-average}, not the entropy of the mixture.

\section{The physical current and its Kac--Moody level}
\label{sec:current}
\subsection{From a symmetry to the full affine algebra}
A conserved internal current generates rotations of fields, whereas the
stress tensor generates changes of coordinates. Their algebras therefore
contain different information. Begin with generators $t^a$ of a compact
Lie algebra $\mathfrak g$, with
$[t^a,t^b]=\ii f^{ab}{}_ct^c$ and invariant inner product $\kappa^{ab}$.
The indices label symmetry generators, not lattice orbitals. The
holomorphic currents have the singular product
\begin{equation}
 J^a(z)J^b(u)\sim
 \frac{k\kappa^{ab}}{(z-u)^2}
 +\frac{\ii f^{ab}{}_cJ^c(u)}{z-u},\qquad
 J_n^a=\oint\frac{\dd z}{2\pi\ii}\,z^nJ^a(z).
 \label{eq:nonabelian-ope}
\end{equation}
The simple pole reproduces the original Lie algebra; the double pole
encodes its central extension. Here and below CFT fields and modes use
the standard unhatted notation, although they are quantum operators;
their central identity is still $\One$.

To see explicitly how the local product gives a commutator, subtract the
two radial orderings of the contours. The remaining inner contour
surrounds $u$ and selects the singular part:
\begin{align}
 [J_m^a,J_n^b]
 &=\oint\frac{\dd u}{2\pi\ii}\,u^n
   \oint_u\frac{\dd z}{2\pi\ii}\,z^m
     \left[\frac{k\kappa^{ab}\One}{(z-u)^2}
          +\frac{\ii f^{ab}{}_cJ^c(u)}{z-u}\right]\nonumber\\
 &=\ii f^{ab}{}_cJ^c_{m+n}
       +k m\kappa^{ab}\delta_{m+n,0}\One.
 \label{eq:full-km}
\end{align}
The double pole differentiates $z^m$; the final contour is nonzero only
when $m+n=0$. In particular, the zero modes satisfy the ordinary global
symmetry algebra with no central term. The nonzero modes describe its
spatially varying current fluctuations. This is the affine Kac--Moody
algebra in a representation of level $k$
\cite{Ginsparg1989}, Sec.~9.1, Eqs.~(9.1)--(9.2).

For completeness, before selecting a representation the central number
is replaced by a generator $\mathsf K$. The extended loop algebra also
includes a grading derivation $d$:
\begin{align}
 [J_m^a,J_n^b]&=\ii f^{ab}{}_cJ^c_{m+n}
          +m\kappa^{ab}\delta_{m+n,0}\mathsf K,\nonumber\\
 [\mathsf K,J_n^a]&=0,\qquad [d,\mathsf K]=0,
       \qquad [d,J_n^a]=nJ_n^a.
 \label{eq:affine-extension}
\end{align}
In a fixed-level irreducible representation $\mathsf K=k\One$.
The derivation records the mode number; it is not an additional physical
charge. Since $[L_0,J_n^a]=-nJ_n^a$, it may be represented by $-L_0$
up to a commuting constant. These definitions distinguish the abstract
affine algebra from its realization in a particular CFT
\cite{HernandezAffine}, Definitions~3.3 and~3.5 and Sec.~4.2.

An antiholomorphic sector has a separate copy of
Eq.~\eqref{eq:full-km}, and the two copies commute. Using $R,L$ for
physical propagation directions, their levels are $k_R,k_L$; conformal
symmetry alone does not require them to be equal. Each current is a
weight-one primary, so the combined conformal and current algebra also
contains
\begin{equation}
 [L_m,J_n^a]=-nJ_{m+n}^a,
 \qquad [L_m,\bar J_n^a]=[J_m^a,\bar J_n^b]=0,
 \label{eq:virasoro-current}
\end{equation}
with the corresponding barred relation. Together with the two Virasoro
algebras, this specifies the symmetry structure used here. The first
commutator follows from the weight-one version of
Eq.~\eqref{eq:T-primary} \cite{Ginsparg1989}, Eq.~(9.8).

\subsection{The Abelian current, its states, and physical normalization}
For $U(1)$ there is one generator and $f^{ab}{}_c=0$. Set its metric to
one, obtaining
\begin{equation}
 J(z)J(0)\sim\frac{k}{z^2},\qquad
 J(z)=\sum_{n\in\mathbb Z}J_nz^{-n-1},
 \label{eq:current-ope}
\end{equation}
and
\begin{equation}
 [J_m,J_n]=km\,\delta_{m+n,0}\One,
 \qquad J_n^\dagger=J_{-n}.
 \label{eq:km}
\end{equation}
Thus the Abelian algebra is a set of harmonic oscillators together with
the conserved charge $J_0$. A current highest-weight state is defined by
\begin{equation}
 J_{n>0}\ket Q=0,\qquad J_0\ket Q=Q\ket Q.
 \label{eq:current-highest-weight}
\end{equation}
Applying $J_{-n}$ creates a density-wave excitation without changing
$Q$. For a normalized highest-weight state,
$\|J_{-n}\ket Q\|^2=\bra Q[J_n,J_{-n}]\ket Q=kn$ for $n>0$.
Unitarity requires $k>0$ for a nontrivial current. It does not, from this
oscillator calculation alone, require an integer $k$. The word
``level'' refers to the central-extension coefficient, not an occupation
number or an individual energy level.

The physical charge normalization is fixed by its action on local
operators. If $\mathcal O_Q$ carries charge $Q$, then
\begin{equation}
 J(z)\mathcal O_Q(u)\sim\frac{Q\mathcal O_Q(u)}{z-u},
 \qquad [J_0,\mathcal O_Q]=Q\mathcal O_Q.
 \label{eq:charge-ward}
\end{equation}
This contour Ward identity connects a continuum normalization to the
microscopic number operator. A physical annihilation field has charge
$-1$ and its adjoint has charge $+1$.

For a charge-one chiral complex fermion with propagator
$\langle\psi(z)\psi^\dagger(0)\rangle=1/z$, the normal-ordered bilinear
current $J=:\!\psi^\dagger\psi\!:$ has double contraction $1/z^2$.
Thus $k=1$. Normal ordering removes the coincident vacuum contraction.
The plane fermion here has a unit-normalized propagator; the conversion
to the physical density retains a factor $1/(2\pi)$ in
Eq.~\eqref{eq:physical-density}. This is compatible with the
$1/\sqrt{2\pi a}$ in Eq.~\eqref{eq:user-bosonization}, not a change
of the physical unit charge.
The normalization is physical, not just notation: replacing $J$ by
$\eta J$ changes $k$ to $\eta^2 k$ and changes the charge generator.
Throughout, the microscopic number operator fixes the unit charge.
For compact simple non-Abelian symmetry, positive-energy unitary
integrable representations require nonnegative integer level in the
basic normalization of the invariant form
\cite{Ginsparg1989}, Sec.~9.3. This statement must not be transplanted
without qualification to a physical $U(1)$ component. Compactification,
the allowed charge lattice, and the combination of chiralities matter.
Indeed, a nonchiral compact boson has continuously variable equal
left/right physical current coefficients as its radius is varied
\cite{BenjaminOoguriShaoWang2020}, Sec.~2.3. An integer coefficient in a
free-fermion example is therefore neither a general theorem about
Abelian currents nor a diagnostic of chirality.

\subsection{Density correlator on a circle}
A current is a primary of holomorphic weight one, so under
Eq.~\eqref{eq:cylinder-map},
\begin{equation}
 J_{\mathrm{cyl}}(w)=\frac{2\pi}{L}zJ(z),\qquad
 \hat\rho_w(y)=\frac{1}{2\pi}J_{\mathrm{cyl}}(\ii y).
 \label{eq:physical-density}
\end{equation}
The subscript $w$ labels the wall branch; the argument of
$J_{\mathrm{cyl}}$ is the Euclidean coordinate.
With real-time continuation $\tau=\ii t$, the choice $\tau+\ii y$
is the left-moving orientation: a holomorphic field depends on $t+y$.
The right mover of Eq.~\eqref{eq:chiral-fermion} instead uses
$\tau-\ii y$. All equal-time variances below are the same for either
orientation; their mode-to-density Fourier signs are opposite.
The factor $1/(2\pi)$ makes
$\int_0^L\dd y\,\hat\rho_w(y)=J_0$ for the fluctuating charge relative
to the reference background. At separated points,
\begin{equation}
 \langle J_{\mathrm{cyl}}(w)J_{\mathrm{cyl}}(0)\rangle_c
 =k\left(\frac{2\pi}{L}\right)^2\frac{z}{(z-1)^2}.
\end{equation}
At equal time, $z=\e^{2\pi\ii y/L}$ and
$z/(z-1)^2=-1/[4\sin^2(\pi y/L)]$. Therefore
\begin{equation}
 C_w(y)\equiv\langle\hat\rho_w(y)\hat\rho_w(0)\rangle_c
 =-\frac{k_w}{4\pi^2}\left(\frac{\pi}{L}\right)^2
    \csc^2\left(\frac{\pi y}{L}\right),\qquad y\ne0 .
 \label{eq:density-correlator}
\end{equation}
The antiholomorphic branch gives the same equal-time expression. Its
orientation does not change the sign of a charge variance.
The negative separated-point correlation is compatible with a positive
variance: a positive short-distance contact contribution must also be
included. Equation~\eqref{eq:density-correlator} is not an unregulated
formula at $y=0$.

\subsection{Integrating the correlator}
Let $\hat Q_{[0,\ell]}=\int_0^\ell\dd y\,\hat\rho_w(y)$, with a
short-distance regulator. Translation invariance and the symmetry of the
equal-time connected correlator imply
\begin{equation}
 F_w(\ell)=\int_0^\ell\dd y\int_0^\ell\dd y'\,C_w(y-y'),\qquad
 F_w''(\ell)=2C_w(\ell)
 \label{eq:variance-derivatives}
\end{equation}
away from the regulated endpoint region. Since
$\dd^2\log\sin(\pi\ell/L)/\dd\ell^2
=-(\pi/L)^2\csc^2(\pi\ell/L)$, integration gives
\begin{equation}
 F_w(\ell)=\frac{k_w}{2\pi^2}
       \log\sin\left(\frac{\pi\ell}{L}\right)+A_w\ell+B_w.
 \label{eq:integrated-current}
\end{equation}
The separated-point correlator alone leaves the linear term undetermined.
In a definite-total-charge state, $Q_{\bar A}=Q_{\mathrm{tot}}-Q_A$ gives
$F(\ell)=F(L-\ell)$, after invoking translation invariance. This symmetry
sets $A_w=0$ for an independent fixed-charge branch. The contact term and
cutoff then enter the additive constant:
\begin{equation}
 F_w(\ell)=\frac{k_w}{2\pi^2}X_L(\ell)+b_{F,w}+\cdots,\qquad
 a_F=\frac{k_\Sigma}{2\pi^2}.
 \label{eq:charge-slope}
\end{equation}
Here $k_\Sigma=\sum_w k_w$ counts independent branches. If only the
combined system has definite total charge, the argument applies to its
full variance; separate wall variances require additional independence
and zero-mode assumptions. Grand-canonical mixing, fluctuating zero modes,
or finite-temperature states can supply other terms. Dropping the linear
term is consequently a physical condition, not just a choice of integration
constant.

The CFT expression is valid for $a\ll\ell,L-\ell$. It is not to be
extrapolated to $\ell=0$ or $L$, where the exact lattice variance is zero
and the logarithmic continuum approximation is invalid.

\subsection{A regulated mode sum: where the contact term went}
The oscillator algebra supplies a complementary derivation that keeps
positivity and the endpoints visible. Equation~\eqref{eq:physical-density}
gives $\hat\rho_w(y)=L^{-1}\sum_nJ_n\e^{-2\pi\ii ny/L}$.
In a highest-weight state of fixed charge, the zero mode has no variance
and $\langle J_nJ_{-n}\rangle_c=kn$ for $n>0$. Integrating the Fourier
modes over $[0,\ell]$ and regulating them by
$r^n=\exp(-2\pi a n/L)$ gives
\begin{align}
 F_{w,a}(\ell)
 &=\frac{k_w}{\pi^2}\sum_{n=1}^\infty
             \frac{r^n}{n}\sin^2\frac{\pi n\ell}{L}\nonumber\\
 &=\frac{k_w}{4\pi^2}
       \log\left[1+\frac{4r\sin^2(\pi\ell/L)}{(1-r)^2}\right].
 \label{eq:regulated-current-sum}
\end{align}
To obtain the second line, write $2\sin^2\theta=1-\cos2\theta$ and
use $\sum_{n>0}z^n/n=-\log(1-z)$. This expression is nonnegative,
symmetric under $\ell\leftrightarrow L-\ell$, and vanishes at both
endpoints. For $a\ll\ell,L-\ell$, $1-r\simeq2\pi a/L$ reproduces
Eq.~\eqref{eq:charge-slope}. It is an explicit continuum regulator, not
a universal finite-lattice formula. The short-distance contribution
needed to compensate the negative separated-point correlator is already
contained in the regulated sum. A symmetrized density covariance may be
used throughout this variance calculation; the antisymmetric contact
commutator integrates to zero against the same interval on both sides.

The zero-mode assumption can now be read off rather than hidden in an
integration constant. An incoherent mixture of highest-weight sectors
with fluctuating $J_0$ contributes
\begin{equation}
 F_{w,a}^{\mathrm{mix}}(\ell)
 =F_{w,a}(\ell)+\left(\frac{\ell}{L}\right)^2\Var(J_0).
 \label{eq:zero-mode-variance}
\end{equation}
Excited oscillator populations give further terms. Fixed charge alone
does not make an arbitrary state a current vacuum; both conditions entered
the derivation. Equation~\eqref{eq:zero-mode-variance} is also a continuum
example of the conditional-versus-mixture distinction in
Eq.~\eqref{eq:total-variance}.

\section{Why central charge and current level are distinct}
\label{sec:sugawara}
The central charge normalizes the stress-tensor OPE,
\begin{equation}
 T(z)T(0)\sim\frac{c/2}{z^4}
          +\frac{2T(0)}{z^2}+\frac{\partial T(0)}{z}.
 \label{eq:stress-ope}
\end{equation}
The current level instead normalizes Eq.~\eqref{eq:current-ope}.
Their relation can be understood directly for an Abelian current.

\subsection{Sugawara construction and charged descendants}
Define its Sugawara stress tensor by
\begin{equation}
 T_J(z)=\frac{1}{2k}:\!J(z)J(z)\!: .
 \label{eq:sugawara}
\end{equation}
The two possible contractions with $J(0)$ give
$T_J(z)J(0)\sim J(0)/z^2+\partial J(0)/z$, so the current has
weight one under this stress tensor. In $T_J(z)T_J(0)$ there are two
double contractions and four single contractions. Their coefficients are
\begin{equation}
 \frac{2k^2}{(2k)^2z^4}=\frac{1}{2z^4},\qquad
 \frac{4k}{(2k)^2z^2}:\!J(z)J(0)\!:
 =\frac{2T_J(0)}{z^2}+\frac{\partial T_J(0)}{z}+\cdots .
 \label{eq:sugawara-contractions}
\end{equation}
Comparison with Eq.~\eqref{eq:stress-ope} gives $c_J=1$, independent of
$k>0$. This is the Abelian Sugawara construction
\cite{DiFrancescoMathieuSenechal1997}, and already shows why $c=k$
is not a general identity.

The corresponding modes display what this stress tensor counts:
\begin{equation}
 L_n^{(J)}=\frac{1}{2k}\sum_{p\in\mathbb Z}:J_{n-p}J_p:,\qquad
 L_0^{(J)}=\frac{J_0^2}{2k}+\frac1k\sum_{p>0}J_{-p}J_p .
 \label{eq:sugawara-modes}
\end{equation}
With $b_p=J_p/\sqrt{kp}$ for $p>0$, the second term is
$\sum_{p>0}p\,b_p^\dagger b_p$. It counts density-wave excitations above
the charged highest-weight state; the first term is its charge cost.
Thus the eigenvalue of $L_0^{(J)}$ on a descendant is
$Q^2/(2k)+\sum_{p>0}pN_p$. The cylinder energy of this chiral sector is
$(2\pi v/L)(L_0^{(J)}-1/24)$, with the last term the Casimir shift.
This makes the current representation an explicit spectrum rather than
only a commutator.

For a compact simple non-Abelian algebra, now choose $\kappa$ to be the
basic invariant metric (long roots have squared length two), and define
the level $k$ relative to that choice. The analogous result is
\begin{equation}
 T_{\mathfrak g}
   =\frac{\kappa_{ab}:J^aJ^b:}{2(k+h^\vee)},\qquad
 c_{\mathfrak g}=\frac{k\,\dim\mathfrak g}{k+h^\vee}.
 \label{eq:nonabelian-sugawara}
\end{equation}
Here $\kappa_{ab}$ is the inverse metric and $h^\vee$ is the dual Coxeter
number. The shift
$h^\vee$ is the extra contribution of the non-Abelian simple pole when
normal ordering the currents. It is absent for $U(1)$, recovering
Eq.~\eqref{eq:sugawara}. Specifying the metric matters: rescaling it
changes the number called $k$. These formulas and their normal-ordering
derivation are given in Ref.~\cite{Ginsparg1989}, Sec.~9.2.

\subsection{Neutral sectors and the limits of coefficient matching}
When the full theory contains a commuting neutral sector,
$T=T_J+T_{\mathrm{neutral}}$ and
$c=1+c_{\mathrm{neutral}}$, while the physical $U(1)$ current still has
level $k$. Conversely, rescaling the physical charge of a single species
changes $k$ quadratically but leaves its central charge unchanged.
A charged primary with
$J(z)\mathcal O_Q(0)\sim Q\mathcal O_Q(0)/z$
has charged-sector weight $h_Q=Q^2/(2k)$; a neutral factor can add to its
total weight. This follows either by the double contraction in
Eq.~\eqref{eq:sugawara} or by acting with its zero mode on a
current-algebra highest-weight state.

For one charge-one free complex fermion, the current has $k=1$ and its
stress tensor is exhausted by $T_J$, giving $c=1$. Independent copies add
both coefficients. This special theory predicts
\begin{equation}
 a_{S_q}=\frac{\pi^2}{6}\left(1+\frac1q\right)a_F,\qquad
 a_{S_q}-\frac{\pi^2}{6}\left(1+\frac1q\right)a_F
 =\frac{c_\Sigma-k_\Sigma}{12}\left(1+\frac1q\right).
 \label{eq:slope-relation}
\end{equation}
These relations concern leading logarithmic coefficients. Equality of two
coefficients by itself does not reconstruct an operator spectrum or
uniquely identify a CFT; it is a consequence of the specified free-fermion
content, not its converse.

\section{A common susceptibility spectrum}
\label{sec:susceptibility}
\subsection{Exact spectral relations}
There is also a microscopic reason why the entropy and variance are closely
related in a Gaussian state. Define
\begin{equation}
 \mathcal V_A=G_A(\Id_A-G_A),\qquad
 v_a=\nu_a(1-\nu_a)\in[0,1/4].
 \label{eq:V}
\end{equation}
Because each entropy kernel is symmetric under $\nu\leftrightarrow1-\nu$,
it is a function of $v$:
\begin{equation}
 F_A=\tr\mathcal V_A,\qquad
 S_q(A)=\tr g_q(\mathcal V_A),\qquad
 g_q(v)=\frac{\log[\nu(v)^q+(1-\nu(v))^q]}{1-q},\quad
 \nu(v)=\frac{1+\sqrt{1-4v}}{2}.
 \label{eq:V-entropies}
\end{equation}
The other root gives the same entropy. For example,
\begin{equation}
 S_2=-\sum_a\log(1-2v_a),\qquad
 S_3=-\frac12\sum_a\log(1-3v_a).
 \label{eq:q23}
\end{equation}
Their small-$v$ expansions begin with $2v$ and $3v/2$, respectively,
whereas $F$ weights every $v$ linearly. Thus the full susceptibility
spectrum fixes all Gaussian R\'enyi entropies, but its first moment $F$
does not. The spectra $\{v_a\}=\{1/4,0\}$ and $\{1/8,1/8\}$ have equal
variance and different $S_2$. The map to $v$ also loses whether each
occupation lies above or below $1/2$, so it cannot determine the mean
charge or odd cumulants.

For a globally pure Slater determinant, $G^2=G$. Taking the $AA$ block
of this identity gives
\begin{equation}
 \mathcal V_A=G_{A\bar A}G_{\bar A A},\qquad
 F_A=\|G_{A\bar A}\|_F^2.
 \label{eq:pure-cross-block}
\end{equation}
The $v_a$ are the squared singular values of the cross-boundary block.
For a mixed Gaussian state the exact replacement is
\begin{equation}
 \mathcal V_A=G_{A\bar A}G_{\bar A A}+(G-G^2)_{AA}.
 \label{eq:mixed-cross-block}
\end{equation}
The positive second term records intrinsic mixedness, making explicit why
the same variance need not represent entanglement in a mixed state.

\subsection{Modular-energy kernels and their integrals}
Use $\epsilon=\log[(1-\nu)/\nu]$ to write
\begin{align}
 f(\epsilon)&=\nu(1-\nu)=\frac{1}{4\cosh^2(\epsilon/2)},\\
 s_q(\epsilon)&=
 \frac{\log[2\cosh(q\epsilon/2)]
                -q\log[2\cosh(\epsilon/2)]}{1-q},\\
 s_1(\epsilon)&=\log[2\cosh(\epsilon/2)]
                   -\frac{\epsilon}{2}\tanh(\epsilon/2).
 \label{eq:kernels}
\end{align}
With the unnormalized spectral density
$\varrho_A(\epsilon)=\sum_a\delta(\epsilon-\epsilon_a)$,
\begin{equation}
 F_A=\int\dd\epsilon\,\varrho_A(\epsilon)f(\epsilon),\qquad
 S_q(A)=\int\dd\epsilon\,\varrho_A(\epsilon)s_q(\epsilon).
 \label{eq:kernel-observables}
\end{equation}
It is the density including the number of modes, not a unit-area histogram,
that enters these equations.

The variance kernel integrates to one because
$f(\epsilon)=-\dd\nu/\dd\epsilon$ and $\nu$ falls from one to zero.
For the entropy kernel, evenness and $\epsilon>0$ give
\begin{equation}
 s_q(\epsilon)=
 \frac{\log(1+\e^{-q\epsilon})-q\log(1+\e^{-\epsilon})}{1-q}.
\end{equation}
For $b>0$, expanding the logarithm and integrating term by term yields
\begin{equation}
 \int_0^\infty\dd\epsilon\,\log(1+\e^{-b\epsilon})
 =\frac1b\sum_{n=1}^\infty\frac{(-1)^{n+1}}{n^2}
 =\frac{\pi^2}{12b}.
\end{equation}
It follows that
\begin{equation}
 \int_{-\infty}^\infty f(\epsilon)\dd\epsilon=1,\qquad
 \int_{-\infty}^\infty s_q(\epsilon)\dd\epsilon
 =\frac{\pi^2}{6}\frac{q^{-1}-q}{1-q}
 =\frac{\pi^2}{6}\left(1+\frac1q\right).
 \label{eq:kernel-integrals}
\end{equation}
The continuous $q=1$ value is $\pi^2/3$.

Suppose the coefficient of the growing spectral density is approximately
energy independent over the kernel-supporting energies,
$\varrho_A(\epsilon)=\gamma X_L(\ell)+\varrho_{\mathrm{sub}}(\epsilon,\ell)$,
where the remainder contributes $o(X_L)$ to both integrals. Then
Eq.~\eqref{eq:kernel-integrals} reproduces Eq.~\eqref{eq:slope-relation}.
This is an additional spectral assumption, not a theorem for every
Gaussian state. More generally an energy-dependent coefficient
$\gamma(\epsilon)$ is weighted differently by $s_q$ and $f$.
A finite density near zero or a logarithmically growing count in one
energy window does not establish an energy-independent coefficient.
Nor are the two kernels proportional pointwise. The derivation explains
a leading-coefficient relation without asserting $S_q(\ell)\propto F(\ell)$
at every finite interval length.

\subsection{Why a homogeneous critical chain can have a flat coefficient}
The spectral assumption has a familiar controlled realization. For a
long interval of a homogeneous critical hopping chain, the active
single-particle modular levels are asymptotically equally spaced, with
spacing of order $\pi^2/\log L$ on each side of zero. Thus the signed
modular-energy density has leading value $\log L/\pi^2$ in that
two-endpoint geometry \cite{PeschelEisler2009}, Sec.~4.1, Eq.~(28).
Inserting this value in Eq.~\eqref{eq:kernel-integrals} gives
$F\simeq\log L/\pi^2$ and
$S_q\simeq(1+q^{-1})\log L/6$, the nonchiral Dirac coefficients.
The logarithm counts the number of partially occupied modes on which
the kernels have appreciable weight. Neither all lattice modes nor just
the mode closest to $\nu=1/2$ determines that coefficient. A defect makes
this explanation more informative by changing the energy coordinate in
which the levels are approximately uniform.

\section{A solvable interface defect: logarithms without homogeneous coefficients}
\label{sec:defect}
Eisler and Peschel's solution \cite{EislerPeschel2010} gives a concrete
test of the preceding reasoning. Their principal example is a critical
transverse Ising (TI) chain of $2L_h$ sites, divided into two halves at a
modified bond. The number-conserving XX hopping chain is obtained from
two Ising sectors and has twice the leading entropy. This is an
equilibrium defect \emph{at an entanglement interface}, not a
demonstration that an adaptive domain wall or a hard-wall measurement
support realizes the same scatterer.

\subsection{The spectral map and the logarithmic mode count}
Write $t_{\mathrm d}>0$ for the defect bond relative to the homogeneous
bond, and introduce
\begin{equation}
 s=\frac{2t_{\mathrm d}}{1+t_{\mathrm d}^2},\qquad
 0\leq s\leq1,\qquad \mathcal T_{\mathrm{def}}=s^2.
 \label{eq:defect-transmission}
\end{equation}
The last equality is the Fermi-energy transmission probability:
it follows from the scattering phases in
Ref.~\cite{EislerPeschel2010}, Eqs.~(31)--(32). Strong and weak bonds
$t_{\mathrm d}$ and $1/t_{\mathrm d}$ have the same $s$.
The spectral relation in their Eqs.~(17)--(20) is
\begin{equation}
 \cosh\Omega_s(\xi)=\frac{\cosh\xi}{s},
 \qquad \xi\geq0,\qquad
 |\epsilon|=2\Omega_s(\xi).
 \label{eq:defect-energy-map}
\end{equation}
We renamed their homogeneous half-energy $\varepsilon$ as $\xi$.
Their transfer matrix supplies $\Omega_s$; the reduced density matrix
requires its square, explaining the factor of two in the actual modular
energy $\epsilon$. Missing this factor would change every entropy
integral.

The geometry behind their Eq.~(21) is also instructive. Mapping an
annulus to a strip by a logarithm produces a strip width proportional to
$\log L_h$. In the notation of that paper its aspect ratio obeys
$2M/N=2\pi/\log L_h$, and hence
\begin{equation}
 \sum_{\text{TI modes}}\longrightarrow
 \frac{N}{M\pi}\int_0^\infty\dd\xi
 =\frac{\log L_h}{\pi^2}\int_0^\infty\dd\xi .
 \label{eq:defect-reference-measure}
\end{equation}
The measure is flat in the \emph{reference half-energy} $\xi$, not in
the physical defect energy $2\Omega_s(\xi)$. The logarithmic prefactor
survives because the mapped strip becomes wider with size even though
the dispersion in that strip has changed. Boundary quantization shifts
do not affect this leading continuum density.

Using the even modular kernels from Eq.~\eqref{eq:kernels}, define
\begin{equation}
 \mathcal I_q(s)=\int_0^\infty\dd\xi\,
                     s_q\!\left(2\Omega_s(\xi)\right).
 \label{eq:defect-renyi-integral}
\end{equation}
Then, for this single-cut half-chain geometry,
\begin{equation}
 S_q^{\mathrm{TI}}=\frac{\mathcal I_q(s)}{\pi^2}\log L_h+O(1),
 \qquad
 S_q^{\mathrm{XX}}=\frac{2\mathcal I_q(s)}{\pi^2}\log L_h+O(1).
 \label{eq:defect-entropies}
\end{equation}
The XX factor of two already includes its two Ising sectors; one must
not double it again when using a signed modular-energy density.
For a generic bond defect these relations describe the leading scaling
limit. The follow-up \cite{PeschelEisler2012}, Sec.~II, constructs
special scale-free defects for which the spectral identities hold more
directly at finite size. They are not identities for arbitrary
inhomogeneous Gaussian matrices.

\subsection{The entropy integral and effective central charge}
At $q=1$, Eq.~\eqref{eq:defect-renyi-integral} is precisely the
$I=I_1+I_2$ of Ref.~\cite{EislerPeschel2010}, Eqs.~(22)--(23):
\begin{equation}
 \mathcal I_1(s)=\int_0^\infty\dd\xi\,
 \left[\log(1+\e^{-2\Omega_s})
              +\frac{2\Omega_s}{1+\e^{2\Omega_s}}\right].
 \label{eq:defect-von-neumann-integral}
\end{equation}
Their crucial simplification is to differentiate with respect to $s$.
The two needed derivatives are
\begin{equation}
 \partial_s\Omega_s=-\frac{\coth\Omega_s}{s},\qquad
 \frac{\dd}{\dd\Omega}s_1(2\Omega)=-\Omega\operatorname{sech}^2\Omega,
 \qquad
 \mathcal I_1'(s)=\frac1s\int_0^\infty
       \frac{\Omega_s\,\dd\xi}{\sinh\Omega_s\cosh\Omega_s}.
 \label{eq:defect-first-derivative}
\end{equation}
Differentiating once more and substituting $x=\cosh\Omega_s=\cosh\xi/s$
produces the explicit intermediate integral in their Eq.~(24):
\begin{equation}
 \mathcal I_1''(s)=\frac{1}{s^2}\int_{1/s}^{\infty}
 \frac{\dd x}{\sqrt{x^2-s^{-2}}}
 \left[\frac{x\operatorname{arcosh}x}{(x^2-1)^{3/2}}
                         -\frac{1}{x^2-1}\right].
 \label{eq:defect-second-derivative-integral}
\end{equation}
The bracket equals
$-\dd[\operatorname{arcosh}x/\sqrt{x^2-1}]/\dd x$.
Evaluating the convergent integral, with the combined endpoint terms
kept together when performing integrations by parts, gives their
Eq.~(25):
\begin{equation}
 \mathcal I_1''(s)=-\frac{\log s}{1-s^2},\qquad
 \mathcal I_1(0)=\mathcal I_1'(0)=0.
 \label{eq:defect-dilog-ode}
\end{equation}
The boundary conditions fix both integration constants. Integrating
twice, or differentiating the expression below to verify them, yields
\begin{align}
 \mathcal I_1(s)=-\frac12\bigl\{&
 [(1+s)\log(1+s)+(1-s)\log(1-s)]\log s\nonumber\\
 &+(1+s)\operatorname{Li}_2(-s)
                    +(1-s)\operatorname{Li}_2(s)\bigr\},\qquad
 \operatorname{Li}_2(z)=-\int_0^z\frac{\log(1-x)}{x}\dd x .
 \label{eq:defect-dilog}
\end{align}
Limits at $s=0,1$ are taken continuously. In particular
$\mathcal I_1(1)=\pi^2/12$, so the transparent TI half-chain has
$S_1=\log L_h/12$ and the transparent XX half-chain has
$S_1=\log L_h/6$, as their respective bulk central charges $1/2$ and
$1$ require for one entangling cut.

For the XX problem it is convenient to write
\begin{equation}
 S_1^{\mathrm{XX}}=\frac{c_{\mathrm{eff}}(s)}6\log L_h+O(1),
 \qquad c_{\mathrm{eff}}(s)=\frac{12\mathcal I_1(s)}{\pi^2}.
 \label{eq:defect-ceff}
\end{equation}
This coefficient varies continuously from zero to one. The bulk remains
the same critical Dirac theory: its local stress-tensor algebra still
has $c_R=c_L=1$. What changed is transmission across the entanglement
interface. An effective entropy coefficient is not automatically a
Virasoro central charge, a current level, or a count of chiral branches.

The Rényi dependence is equally explicit. Evaluating
Eq.~\eqref{eq:defect-renyi-integral}, as in
Ref.~\cite{PeschelEisler2012}, Sec.~IV, gives
\begin{align}
 \mathcal I_2(s)&=\arcsin^2(s/\sqrt2),&
 \mathcal I_3(s)&=\frac12\arcsin^2(\sqrt3s/2),\nonumber\\
 \mathcal I_{1/2}(s)&=\frac12\arcsin s\,(\pi-\arcsin s),&
 \mathcal I_\infty(s)&=\frac14\operatorname{Li}_2(s^2).
 \label{eq:defect-renyi-examples}
\end{align}
For example, $s_2(2\Omega)=-\log[1-s^2/(2\cosh^2\xi)]$.
If $A(b)=-\int_0^\infty\log(1-b\,\operatorname{sech}^2\xi)\dd\xi$,
the substitution $u=\tanh\xi$ gives
$A'(b)=\arcsin\sqrt b/\sqrt{b(1-b)}$. Since $A(0)=0$,
$A(b)=\arcsin^2\sqrt b$, proving the $q=2$ formula; the $q=3$ kernel
gives the next one in the same way. The $q=\infty$ kernel is
Eq.~\eqref{eq:single-copy}. Its integral is the first term of the
paper's $I_1+I_2$, not our von Neumann $\mathcal I_1$.

At transparency
$\mathcal I_q(1)=\pi^2(1+q^{-1})/24$, recovering the homogeneous
Rényi relation. Away from it, a single $q$-independent replacement of
$c$ in the homogeneous formula does not in general reproduce all
Rényi indices. For example,
\begin{equation}
 \mathcal I_1(s)=s^2\left(\frac34-\frac12\log s\right)
                   +O(s^4|\log s|),\qquad
 \mathcal I_q(s)\sim\frac{q}{4(q-1)}s^2
                   \quad(q>1,\ s\to0).
 \label{eq:defect-small-s}
\end{equation}
The weak-transmission and $q\to1$ limits are nonuniform.
The $q=0$ rank entropy requires a separate order-of-limits analysis.

\subsection{Susceptibility, charge fluctuations, and the nonflat density}
The same spectral map immediately implies
\begin{equation}
 v_{\mathrm{def}}=\frac{1}{4\cosh^2\Omega_s(\xi)}
       =s^2\frac{1}{4\cosh^2\xi}
       =\mathcal T_{\mathrm{def}}v_{\mathrm{ref}}.
 \label{eq:defect-susceptibility}
\end{equation}
This is the occupation-product relation derived in
Ref.~\cite{PeschelEisler2012}, Sec.~II, Eqs.~(6)--(9).
For the number-conserving XX chain it gives
\begin{equation}
 F^{\mathrm{XX}}=
 \frac{2\log L_h}{\pi^2}\int_0^\infty
                   \frac{s^2\,\dd\xi}{4\cosh^2\xi}+O(1)
 =\frac{\mathcal T_{\mathrm{def}}}{2\pi^2}\log L_h+O(1).
 \label{eq:defect-charge}
\end{equation}
We do not interpret a TI Bogoliubov-mode occupation variance as a
conserved physical $U(1)$ charge variance: the TI model does not have
that number symmetry. In the XX model the entropy is nonlinear in
$v_{\mathrm{def}}$ but the variance is linear, so their leading ratio is
\begin{equation}
 \frac{a_{S_1}^{\mathrm{XX}}}{a_F^{\mathrm{XX}}}
   =\frac{4\mathcal I_1(s)}{s^2}.
 \label{eq:defect-ratio}
\end{equation}
It equals $\pi^2/3$ at $s=1$, not at general transmission. A naive
homogeneous conversion of the suppressed variance coefficient into
a level would therefore measure transmission-weighted information,
not a change of the bulk current algebra.

Finally change variables from $\xi$ to the actual signed XX modular
energy $\epsilon=\pm2\Omega_s(\xi)$. The Jacobian gives, for $0<s<1$,
\begin{align}
 \varrho_s(\epsilon)
 &=\frac{\log L_h}{2\pi^2}
   \frac{s\sinh(|\epsilon|/2)}
        {\sqrt{s^2\cosh^2(\epsilon/2)-1}}\,
   \Theta_{\mathrm H}(|\epsilon|-\epsilon_*)+
      \varrho_{\mathrm{sub}}(\epsilon),\nonumber\\
 \epsilon_*&=2\operatorname{arcosh}(1/s).
 \label{eq:defect-density}
\end{align}
Here $\Theta_{\mathrm H}$ is the Heaviside function. To check the
normalization, use
$\dd\xi/\dd|\epsilon|
=s\sinh(|\epsilon|/2)/[2\sqrt{s^2\cosh^2(\epsilon/2)-1}]$
and count both signs once. The density has a modular gap and an
integrable square-root threshold, becoming flat only at $s=1$.
Nevertheless its overall amplitude still grows as $\log L_h$.
This explains how a gapped entanglement spectrum and a logarithmic
entropy coexist, without gapping the physical critical chain.

The lesson for the adaptive-wall interpretation is a logical one.
Before using a logarithmic coefficient as a CFT datum, one must specify
the state, the entangling geometry, and any transmitting interface.
The solvable defect supplies a counterexample to an inference, not
evidence that the adaptive circuit has that defect spectrum or that its
hard-wall support parameter equals $s$.

\section{What can be inferred about individual walls?}
\label{sec:walls}
\subsection{One-wall and two-wall normalization}
For a full-width strip
$A(y_0,\ell)=[0,N_x)\times[y_0,y_0+\ell)$ on a torus, both interfaces
are intersected. The two endpoints of the interval are already included
in the twist two-point function; one must not count them as an additional
pair of wall branches. Table~\ref{tab:normalization} collects the
normalizations for charge-one complex fermions.

\begin{table}[tb]
\caption{Leading coefficients of $X_L(\ell)$ for independent free branches.
The last three columns are slopes, not the entropies or variances themselves.
Opposite propagation orientations cancel signed central charge but not
these positive coefficients.}
\label{tab:normalization}
\centering
\begin{ruledtabular}
\begin{tabular}{lccccc}
Content & $c_\Sigma$ & $k_\Sigma$ & $a_{S_q}$ & $a_{S_1}$ & $a_F$\\
\hline
One chiral complex branch & $1$ & $1$ & $(1+q^{-1})/12$ & $1/6$ & $1/(2\pi^2)$\\
Two opposite chiral walls & $2$ & $2$ & $(1+q^{-1})/6$ & $1/3$ & $1/\pi^2$\\
One nonchiral Dirac channel & $2$ & $2$ & $(1+q^{-1})/6$ & $1/3$ & $1/\pi^2$
\end{tabular}
\end{ruledtabular}
\end{table}

For two identical single-branch walls, extracting the per-wall quantities
from full-strip slopes gives
\begin{equation}
 c_w=\frac{6a_{S_q}}{1+1/q},\qquad k_w=\pi^2a_F .
 \label{eq:two-wall-estimators}
\end{equation}
For a genuinely isolated single branch the conversion instead is
$c_w=12a_{S_q}/(1+1/q)$ and $k_w=2\pi^2a_F$.
At $q=1$, these are $3a_{S_1}$ for either wall of an identical pair and
$6a_{S_1}$ for a single branch. In the usual nonchiral convention the
Dirac theory is called $c=1$ because $c_L=c_R=1$; our explicit sum for that
same theory is $c_\Sigma=2$.

For a particular wall,
\begin{equation}
 S_1^{(w)}(\ell)=
 \frac{c_L^{(w)}+c_R^{(w)}}{6}X_L(\ell)+\cdots,\qquad
 c_-^{(w)}=c_R^{(w)}-c_L^{(w)} .
 \label{eq:chiral-difference}
\end{equation}
Only if the wall is purely chiral does its entropy slope determine
$|c_-^{(w)}|/6$. In a unitary theory, the general statement is merely
$|c_-^{(w)}|\leq c_L^{(w)}+c_R^{(w)}$.
Static entropy and charge fluctuations supply neither a propagation
direction nor a proof that one chirality is absent. The explicit
left/right replica weights in Ref.~\cite{CastroDetournayIqbalPerlmutter2014},
Sec.~II, make this distinction transparent.

\subsection{Where chirality appears, and why the two-wall total cancels}
Fix one common positive $y$-direction for the whole strip. For two
oppositely propagating charge-one walls, the possible assignment is
\begin{equation}
 (c_R^{(x_R)},c_L^{(x_R)})=(1,0),\qquad
 (c_R^{(x_L)},c_L^{(x_L)})=(0,1).
 \label{eq:opposite-wall-charges}
\end{equation}
Interchanging the wall labels or reversing $y$ reverses the signed
assignment. Each wall has $|c_-^{(w)}|=1$, but the full pair has
$c_R^{\mathrm{tot}}=c_L^{\mathrm{tot}}=1$ and
$c_-^{\mathrm{tot}}=0$. Its entropy nevertheless has slope
$(c_R^{\mathrm{tot}}+c_L^{\mathrm{tot}})/6=1/3$.
There is no contradiction: the entropy adds two positive contributions,
while net chirality subtracts propagation directions. An alternative
boundary-orientation convention reverses the tangent on the other edge;
one must convert both edges to the same direction before taking a total.

The replica derivation contains this distinction before the equal-time
limit is taken. With weights assigned to physical directions,
\begin{equation}
 h_q^R+h_q^L=\frac{c_R+c_L}{24}(q-q^{-1}),\qquad
 h_q^R-h_q^L=\frac{c_R-c_L}{24}(q-q^{-1}).
 \label{eq:twist-dimension-spin}
\end{equation}
The first combination is the twist scaling dimension. The second is
its signed spin (with the chosen orientation), which governs sensitivity
to relative frame rotations or boosts rather than the ordinary
equal-time chord length. This is how entanglement can carry information
about a gravitational anomaly without the static coefficient itself
being $c_R-c_L$ \cite{CastroDetournayIqbalPerlmutter2014},
Sec.~II.2, Eqs.~(21)--(22), and Sec.~II.3.2.

The same sum-versus-difference distinction occurs for charge. To make
the orientation explicit, expand the densities at equal time as
\begin{equation}
 \rho_R(y)=\frac1L\sum_nJ_n^R\e^{2\pi\ii ny/L},\qquad
 \rho_L(y)=\frac1L\sum_nJ_n^L\e^{-2\pi\ii ny/L}.
 \label{eq:oriented-current-modes}
\end{equation}
Both mode algebras have the positive-level form of Eq.~\eqref{eq:km}.
Since $\partial_y\delta_L(y)=
L^{-1}\sum_n(2\pi\ii n/L)\e^{2\pi\ii ny/L}$, the total physical
density has contact commutator
\begin{equation}
 [\rho_R(y)+\rho_L(y),\rho_R(y')+\rho_L(y')]
 =-\frac{\ii(k_R-k_L)}{2\pi}
                  \partial_y\delta_L(y-y')\One .
 \label{eq:charge-anomaly-bracket}
\end{equation}
Its symmetric, separated-point fluctuations instead add the two levels,
giving $a_F=(k_R+k_L)/(2\pi^2)$. Thus $k_-=k_R-k_L$ measures the
signed charged anomaly, while $k_\Sigma=k_R+k_L$ measures the
fluctuation strength. The difference and its global quantization
conditions are discussed in Ref.~\cite{BenjaminOoguriShaoWang2020},
Secs.~2.1--2.2. An ordinary Dirac fermion has two individually
level-one currents and zero $k_-$; quantized levels do not establish
net chirality. Moreover $k_-$ and $c_-$ remain different anomalies:
neutral modes contribute to the latter without contributing to charge.

Accordingly, a full-strip entropy or variance fit can support the
expected total sector content, but distinguishing separated chiral walls
from a nonchiral channel requires spatial information and an independent
direction-sensitive argument. An integer-looking coefficient alone
does not supply that argument, and the defect example shows an additional
reason not to identify every fitted logarithmic coefficient with a
protected algebraic quantity.

\subsection{Entanglement contours and wall windows}
For a Gaussian reduced state define the cell-resolved contour
\cite{ChenVidal2014}
\begin{equation}
 s_q(\boldsymbol r;A)=
 \sum_\mu\langle\boldsymbol r,\mu|
 h_q(G_A)|\boldsymbol r,\mu\rangle
 =\sum_{\mu,a}|U_{\boldsymbol r\mu,a}|^2h_q(\nu_a),
 \label{eq:contour}
\end{equation}
where $h_q(\nu)$ is the single-mode entropy kernel in
Eq.~\eqref{eq:renyi}. For $q=1$ it is Eq.~\eqref{eq:von-neumann}.
The contour is nonnegative and obeys the exact sum rule
\begin{equation}
 \sum_{\boldsymbol r\in A}s_q(\boldsymbol r;A)=S_q(A).
 \label{eq:contour-sum}
\end{equation}
The diagonal is unchanged by transposing the matrix function, so this
formula is consistent with Eq.~\eqref{eq:transpose}.
Likewise $f(\boldsymbol r;A)=
\sum_\mu\langle\boldsymbol r,\mu|\mathcal V_A|\boldsymbol r,\mu\rangle$
sums to $F_A$.

For a window $W_w$ surrounding one wall, define its assigned contribution
\begin{equation}
 S_q^{[W_w]}(A)=\sum_{\boldsymbol r\in A\cap W_w}s_q(\boldsymbol r;A).
 \label{eq:wall-window}
\end{equation}
This is not generally $S_q(A\cap W_w)$. The former uses the matrix
function of the original $G_A$; the latter first restricts $G$ to a different
subsystem and then applies the nonlinear function. Restriction and taking
the entropy kernel do not commute.

To interpret the coefficient of Eq.~\eqref{eq:wall-window} as the
contribution of an independent wall branch, the window must capture that
branch's leading logarithmic weight; coupling to other walls must be
negligible at the scales in question; and excluded tails must not carry an
unaccounted logarithmic contribution. The contour sum rule alone proves
none of these statements. Contours distribute a total entropy according
to a specified prescription rather than defining a unique local observable.
The interpretation of a finite window therefore needs a localization
argument as well as the exact Gaussian formula.

Under those assumptions, one charge-one chiral wall contributes a von
Neumann slope $a_{S_1}^{[W_w]}=1/6$, or weight
$3a_{S_1}^{[W_w]}=1/2$ in the usual full-strip
normalization. The two opposite walls contribute $1/2+1/2=1$ while their
signed chiral central charges cancel. Without those assumptions, the
window slope remains a well-defined assigned entropy coefficient, not
automatically a central charge.

\section{Scope of the derivations}
The correlation-spectrum expressions, counting determinant, cumulants,
susceptibility identities, and contour closure are exact for the specified
number-conserving Gaussian states. They apply before record averaging.
The twist-field and current-correlator calculations instead assume the
stated infrared CFT and its vacuum-like interval state, with the physical
charge normalization held fixed. Passing from a full strip to independent
wall contributions adds a localization and decoupling assumption.

The continuum actions specify useful comparison theories, not the
adaptive evolution law. In particular a bosonically Gaussian Luttinger
liquid is not necessarily a fermionically Gaussian state. The exactly
solvable defect makes another distinction explicit: a critical bulk can
retain its operator algebras while an interface changes the entropy
coefficient and the modular spectral density. Neither the Luttinger
parameter nor the defect transmission is an additional parameter of the
adaptive protocol unless a separate mapping establishes it.

These distinctions explain both the usefulness and the limits of the
comparison. Entropy and charge fluctuations probe different functions of
one microscopic spectrum and different data of a continuum theory.
Agreement with a charge-one complex-fermion prediction is informative,
but it does not turn a static unsigned observable into a measurement of
chirality or establish a preparation-depth bound.

\clearpage
\bibliographystyle{apsrev4-2}
\bibliography{references}
\end{document}
